\documentclass[11pt,a4paper]{article}

\usepackage[T1]{fontenc}
\usepackage[utf8]{inputenc}
\usepackage[french]{babel}
\usepackage[margin=2.5cm]{geometry}
\usepackage{booktabs}
\usepackage{siunitx}
\usepackage[table]{xcolor}
\usepackage{hyperref}
\usepackage{longtable}
\usepackage{enumitem}
\usepackage{amsmath}
\usepackage{amssymb}

\sisetup{detect-weight=true, detect-family=true}

\hypersetup{
  colorlinks=true,
  linkcolor=blue!50!black,
  urlcolor=blue!50!black,
  citecolor=blue!50!black,
  pdftitle={Rapport d'activite - Campagne PBE vs r2SCAN Al4C3},
  pdfauthor={Pr. John Akapulco}
}

\title{\vspace{-1.5cm}Rapport d'activité\\[2pt]
\large Campagne r2SCAN d'Al$_4$C$_3$ : méthodologie, structures d'entrée et maillage
$\mathbf{k}$, en vue de la comparaison avec le scan PBE}
\author{Pr. John Akapulco \\ \small \href{mailto:computchem86@gmail.com}{computchem86@gmail.com}}
\date{23 août 2026}

\begin{document}
\maketitle
\vspace{-0.5cm}

\begin{abstract}
\noindent
Ce rapport documente le lancement de la campagne r2SCAN d'Al$_4$C$_3$, destinée à être
comparée au scan PBE complet (\texttt{rapport\_activite\_Al4C3\_HP\_scan.tex}, 7 polymorphes,
0--100~GPa). Les paramètres DFT (\texttt{ENCUT}, \texttt{EDIFF}, \texttt{KSPACING}, la grille de
pression, \texttt{ISIF=8}) sont identiques à ceux du scan PBE ; seule la fonctionnelle change
(PBE~$\to$~r2SCAN), pour une comparaison à un seul paramètre variable. Chaque point r2SCAN est
amorcé depuis le \texttt{WAVECAR}/\texttt{CHGCAR} PBE déjà convergé à la même pression (pratique
recommandée par VASP pour les meta-GGA). Ce document liste, pour chacun des sept polymorphes
convergés en PBE, la structure d'entrée (\texttt{POSCAR}, identique à 0 et 50~GPa entre PBE et
r2SCAN) et le maillage de points $\mathbf{k}$ généré par \texttt{KSPACING=0.2} à ces deux
pressions. Il ne contient \emph{pas encore} les résultats r2SCAN eux-mêmes (H(P), classement de
stabilité) : les chaînes de calcul viennent d'être soumises (voir section~\ref{sec:status}) et
ce rapport sera complété d'une comparaison quantitative PBE/r2SCAN une fois les points
convergés.
\end{abstract}

\section{Contexte et objectifs}

Le scan PBE (0--100~GPa, sept polymorphes, \texttt{KSPACING=0.2}) a identifié une séquence de
stabilité Pnma-VII~$\to$~Pbam~$\to$~Pnma avec deux transitions à 26 et 52~GPa (voir rapport
principal, section « Précision des pressions de transition », pour l'incertitude associée à ces
valeurs). L'objectif de la présente campagne est de vérifier que cette séquence -- et en
particulier les deux pressions de transition -- est robuste au choix de fonctionnelle
d'échange-corrélation : r2SCAN (meta-GGA) corrige une partie des biais systématiques connus de
PBE (notamment la sur-estimation des volumes d'équilibre), ce qui peut déplacer les croisements
d'enthalpie entre phases de volumes différents (cf.\ l'analyse de précision du rapport
principal, où $\delta P_t \propto 1/|\Delta V|$).

\section{Méthodologie}

\subsection{Paramètres DFT et redémarrage depuis PBE}

Les paramètres communs (\texttt{PREC}, \texttt{ENCUT=600~eV}, \texttt{EDIFF=1E-5},
\texttt{ISMEAR/SIGMA=0/0.05}, \texttt{KSPACING=0.2}, \texttt{IBRION/ISIF=2/8},
\texttt{NSW=200}, \texttt{EDIFFG=-0.01}, grille de pression 0--100~GPa par pas de 10~GPa,
\texttt{KPAR=8}) sont repris à l'identique du scan PBE. Les seuls ajouts sont les balises
meta-GGA :

\begin{center}
\begin{tabular}{ll}
\toprule
Balise & Rôle \\
\midrule
\texttt{METAGGA = R2scan} & fonctionnelle meta-GGA r2SCAN (implémentation native VASP) \\
\texttt{LASPH = .TRUE.} & corrections PAW asphériques (nécessaires en meta-GGA) \\
\texttt{ADDGRID = .TRUE.} & grille réelle plus fine pour les termes de gradient de $\tau$ \\
\texttt{LMIXTAU = .TRUE.} & mélange de la densité d'énergie cinétique dans le SCF \\
\bottomrule
\end{tabular}
\end{center}

Pour les sept polymorphes dont le scan PBE est intégralement convergé (0--100~GPa), chaque point
r2SCAN est amorcé \emph{directement} depuis le \texttt{CONTCAR}/\texttt{WAVECAR}/\texttt{CHGCAR}
PBE convergé à la même pression (\texttt{ISTART=1}, \texttt{ICHARG=1}) -- c'est la pratique
recommandée par VASP pour la convergence SCF des meta-GGA, et cela élimine le besoin d'une
chaîne de redémarrage entre pressions voisines (chaque point est indépendant). Pour le
polymorphe \textbf{Fd-3m}, dont le scan PBE est encore en cours (voir
section~\ref{sec:status}), seul le point $P=0$ dispose d'un \texttt{WAVECAR} PBE convergé ; la
chaîne r2SCAN de ce polymorphe est donc construite comme une chaîne ascendante unique
0~$\to$~10~$\to \dots \to$~100~GPa, où seul $P=0$ est amorcé depuis PBE et chaque point suivant
redémarre depuis le point r2SCAN précédent (même mécanisme que la chaîne PBE originale de
Fd-3m).

\section{Maillage de points $\mathbf{k}$}

Le maillage \texttt{KSPACING=0.2} génère un nombre de points $\mathbf{k}$ différent pour chaque
polymorphe (la maille et le réseau réciproque diffèrent), et identique entre PBE et r2SCAN à
pression donnée (seule la fonctionnelle change, pas la géométrie du maillage). Le tableau
ci-dessous liste le maillage de Monkhorst-Pack généré ($N_1 \times N_2 \times N_3$), le nombre
total de points $\mathbf{k}$ et le nombre de points irréductibles (après réduction par symétrie)
à 0 et 50~GPa, extraits des \texttt{OUTCAR} PBE convergés.

\begin{center}
\begin{tabular}{lcccc}
\toprule
Polymorphe & $P$ (GPa) & Maillage $N_1{\times}N_2{\times}N_3$ & $k$-points totaux & $k$-points irréductibles \\
\midrule
Pnma-VII           & 0  & $4\times10\times4$ & 160 & 54 \\
Pnma-VII           & 50 & $4\times11\times4$ & 176 & 54 \\
\midrule
Pbam               & 0  & $4\times7\times11$ & 308 & 72 \\
Pbam               & 50 & $4\times7\times11$ & 308 & 72 \\
\midrule
Pnma               & 0  & $4\times11\times4$ & 176 & 54 \\
Pnma               & 50 & $4\times11\times4$ & 176 & 54 \\
\midrule
R-3m               & 0  & $12\times12\times2$ & 288 & 38 \\
R-3m               & 50 & $12\times12\times2$ & 288 & 38 \\
\midrule
Pnma-III           & 0  & $4\times11\times4$ & 176 & 54 \\
Pnma-III           & 50 & $4\times11\times4$ & 176 & 54 \\
\midrule
anti-CaFe$_2$O$_4$ & 0  & $11\times4\times4$ & 176 & 54 \\
anti-CaFe$_2$O$_4$ & 50 & $12\times4\times4$ & 192 & 63 \\
\midrule
I-43d              & 0  & $5\times5\times5$ & 125 & 10 \\
I-43d              & 50 & $6\times6\times6$ & 216 & 20 \\
\bottomrule
\end{tabular}
\end{center}

Le nombre de points $\mathbf{k}$ irréductibles varie fortement d'un polymorphe à l'autre (de 10
pour I-43d à 72 pour Pbam), reflétant à la fois des mailles de tailles différentes (14 à 28
atomes/maille selon le polymorphe) et des symétries de groupe d'espace différentes. I-43d se
distingue par un maillage cubique dense (5$^3$ à 6$^3$) mais très peu de points irréductibles
(groupe d'espace cubique $I\bar43d$, haute symétrie) ; à l'inverse R-3m, malgré 288 points
totaux, n'en réduit qu'à 38 (symétrie rhomboédrique moins favorable pour ce maillage
anisotrope).

\section{Structures d'entrée (POSCAR, PBE relaxé $\to$ seed r2SCAN)}

Les structures ci-dessous sont les \texttt{CONTCAR} PBE convergés à 0 et 50~GPa, utilisés tels
quels comme \texttt{POSCAR} d'entrée des calculs r2SCAN correspondants
(\texttt{hp\_curves\_r2SCAN/Al4C3\_<polymorphe>/P<pression>/POSCAR}).

\subsubsection*{Pnma-VII -- 0~GPa}
\begin{footnotesize}
\begin{verbatim}
Al4C3 Pnma (CIF findsym, P21/n21/m21/a) 
   1.00000000000000     
     9.6447110324068994    0.0000000000000000    0.0000000000000000
     0.0000000000000000    3.3137960964664095    0.0000000000000000
     0.0000000000000000    0.0000000000000000   10.1672966455862142
   Al   C 
    16    12
Direct
  0.2638924603987636  0.2500000000000000  0.2677086017912220
  0.7638924603987637  0.2500000000000000  0.2322913982087779
  0.7361075396012363  0.7500000000000000  0.7322913982087779
  0.2361075396012364  0.7500000000000000  0.7677086017912221
  0.4687336063193806  0.2500000000000000  0.8920766296137922
  0.9687336063193805  0.2500000000000000  0.6079233703862078
  0.5312663936806195  0.7500000000000000  0.1079233703862078
  0.0312663936806194  0.7500000000000000  0.3920766296137922
  0.5078696080686969  0.2500000000000000  0.6025813718867934
  0.0078696080686969  0.2500000000000000  0.8974186281132066
  0.4921303919313030  0.7500000000000000  0.3974186281132066
  0.9921303919313031  0.7500000000000000  0.1025813718867935
  0.2411735741243224  0.2500000000000000  0.5605001914892204
  0.7411735741243224  0.2500000000000000  0.9394998085107796
  0.7588264258756776  0.7500000000000000  0.4394998085107796
  0.2588264258756777  0.7500000000000000  0.0605001914892204
  0.3845072240206697  0.2500000000000000  0.4241631201653526
  0.8845072240206696  0.2500000000000000  0.0758368798346473
  0.6154927759793304  0.7500000000000000  0.5758368798346473
  0.1154927759793303  0.7500000000000000  0.9241631201653527
  0.3904520877040922  0.2500000000000000  0.0833557375444348
  0.8904520877040921  0.2500000000000000  0.4166442624555652
  0.6095479122959079  0.7500000000000000  0.9166442624555650
  0.1095479122959078  0.7500000000000000  0.5833557375444350
  0.3496738716245690  0.2500000000000000  0.7346401141650764
  0.8496738716245689  0.2500000000000000  0.7653598858349236
  0.6503261283754311  0.7500000000000000  0.2653598858349236
  0.1503261283754311  0.7500000000000000  0.2346401141650764
 
  0.00000000E+00  0.00000000E+00  0.00000000E+00
  0.00000000E+00  0.00000000E+00  0.00000000E+00
  0.00000000E+00  0.00000000E+00  0.00000000E+00
  0.00000000E+00  0.00000000E+00  0.00000000E+00
  0.00000000E+00  0.00000000E+00  0.00000000E+00
  0.00000000E+00  0.00000000E+00  0.00000000E+00
  0.00000000E+00  0.00000000E+00  0.00000000E+00
  0.00000000E+00  0.00000000E+00  0.00000000E+00
  0.00000000E+00  0.00000000E+00  0.00000000E+00
  0.00000000E+00  0.00000000E+00  0.00000000E+00
  0.00000000E+00  0.00000000E+00  0.00000000E+00
  0.00000000E+00  0.00000000E+00  0.00000000E+00
  0.00000000E+00  0.00000000E+00  0.00000000E+00
  0.00000000E+00  0.00000000E+00  0.00000000E+00
  0.00000000E+00  0.00000000E+00  0.00000000E+00
  0.00000000E+00  0.00000000E+00  0.00000000E+00
  0.00000000E+00  0.00000000E+00  0.00000000E+00
  0.00000000E+00  0.00000000E+00  0.00000000E+00
  0.00000000E+00  0.00000000E+00  0.00000000E+00
  0.00000000E+00  0.00000000E+00  0.00000000E+00
  0.00000000E+00  0.00000000E+00  0.00000000E+00
  0.00000000E+00  0.00000000E+00  0.00000000E+00
  0.00000000E+00  0.00000000E+00  0.00000000E+00
  0.00000000E+00  0.00000000E+00  0.00000000E+00
  0.00000000E+00  0.00000000E+00  0.00000000E+00
  0.00000000E+00  0.00000000E+00  0.00000000E+00
  0.00000000E+00  0.00000000E+00  0.00000000E+00
  0.00000000E+00  0.00000000E+00  0.00000000E+00
\end{verbatim}
\end{footnotesize}

\subsubsection*{Pnma-VII -- 50~GPa}
\begin{footnotesize}
\begin{verbatim}
Al4C3 Pnma (CIF findsym, P21/n21/m21/a) 
   1.00000000000000     
     8.9923981068448793    0.0000000000000000    0.0000000000000000
     0.0000000000000000    3.0896699387060820    0.0000000000000000
     0.0000000000000000    0.0000000000000000    9.4796390270578463
   Al   C 
    16    12
Direct
  0.2681833644661538  0.2500000000000000  0.2664767518332596
  0.7681833644661538  0.2500000000000000  0.2335232481667404
  0.7318166355338462  0.7500000000000000  0.7335232481667404
  0.2318166355338462  0.7500000000000000  0.7664767518332596
  0.4656894866399242  0.2500000000000000  0.8916382141646366
  0.9656894866399242  0.2500000000000000  0.6083617858353634
  0.5343105133600758  0.7500000000000000  0.1083617858353634
  0.0343105133600758  0.7500000000000000  0.3916382141646366
  0.5060443207492540  0.2500000000000000  0.6048072554787396
  0.0060443207492540  0.2500000000000000  0.8951927445212604
  0.4939556792507460  0.7500000000000000  0.3951927445212604
  0.9939556792507460  0.7500000000000000  0.1048072554787396
  0.2367900995859671  0.2500000000000000  0.5579400534326382
  0.7367900995859671  0.2500000000000000  0.9420599465673618
  0.7632099004140329  0.7500000000000000  0.4420599465673618
  0.2632099004140329  0.7500000000000000  0.0579400534326382
  0.3846996388742596  0.2500000000000000  0.4253344150536265
  0.8846996388742596  0.2500000000000000  0.0746655849463735
  0.6153003611257404  0.7500000000000000  0.5746655849463735
  0.1153003611257404  0.7500000000000000  0.9253344150536265
  0.3944047199639513  0.2500000000000000  0.0877513538253609
  0.8944047199639513  0.2500000000000000  0.4122486461746391
  0.6055952800360487  0.7500000000000000  0.9122486461746391
  0.1055952800360487  0.7500000000000000  0.5877513538253609
  0.3471076641912347  0.2500000000000000  0.7342183592237674
  0.8471076641912347  0.2500000000000000  0.7657816407762326
  0.6528923358087653  0.7500000000000000  0.2657816407762326
  0.1528923358087653  0.7500000000000000  0.2342183592237674
 
  0.00000000E+00  0.00000000E+00  0.00000000E+00
  0.00000000E+00  0.00000000E+00  0.00000000E+00
  0.00000000E+00  0.00000000E+00  0.00000000E+00
  0.00000000E+00  0.00000000E+00  0.00000000E+00
  0.00000000E+00  0.00000000E+00  0.00000000E+00
  0.00000000E+00  0.00000000E+00  0.00000000E+00
  0.00000000E+00  0.00000000E+00  0.00000000E+00
  0.00000000E+00  0.00000000E+00  0.00000000E+00
  0.00000000E+00  0.00000000E+00  0.00000000E+00
  0.00000000E+00  0.00000000E+00  0.00000000E+00
  0.00000000E+00  0.00000000E+00  0.00000000E+00
  0.00000000E+00  0.00000000E+00  0.00000000E+00
  0.00000000E+00  0.00000000E+00  0.00000000E+00
  0.00000000E+00  0.00000000E+00  0.00000000E+00
  0.00000000E+00  0.00000000E+00  0.00000000E+00
  0.00000000E+00  0.00000000E+00  0.00000000E+00
  0.00000000E+00  0.00000000E+00  0.00000000E+00
  0.00000000E+00  0.00000000E+00  0.00000000E+00
  0.00000000E+00  0.00000000E+00  0.00000000E+00
  0.00000000E+00  0.00000000E+00  0.00000000E+00
  0.00000000E+00  0.00000000E+00  0.00000000E+00
  0.00000000E+00  0.00000000E+00  0.00000000E+00
  0.00000000E+00  0.00000000E+00  0.00000000E+00
  0.00000000E+00  0.00000000E+00  0.00000000E+00
  0.00000000E+00  0.00000000E+00  0.00000000E+00
  0.00000000E+00  0.00000000E+00  0.00000000E+00
  0.00000000E+00  0.00000000E+00  0.00000000E+00
  0.00000000E+00  0.00000000E+00  0.00000000E+00
\end{verbatim}
\end{footnotesize}

\subsubsection*{Pbam -- 0~GPa}
\begin{footnotesize}
\begin{verbatim}
Al4C3 Pbam polymorph (findsym-output) - 
   1.00000000000000     
     9.0566422846240506    0.0000000000000000    0.0000000000000000
     0.0000000000000000    5.2293835261696859    0.0000000000000000
     0.0000000000000000    0.0000000000000000    3.0977367561939881
   Al   C 
     8     6
Direct
  0.7846720435396047  0.1531142601818623 -0.0000000000000000
  0.2153279564603954  0.8468857398181378  0.0000000000000000
  0.7153279564603953  0.6531142601818622 -0.0000000000000000
  0.2846720435396045  0.3468857398181377  0.0000000000000000
  0.0250541557709346  0.2542515312388905  0.5000000000000000
  0.9749458682290674  0.7457484687611025  0.5000000000000000
  0.4749458382290650  0.7542515312388975  0.5000000000000000
  0.5250541317709326  0.2457484687611024  0.5000000000000000
  0.0000000000000000 -0.0000000000000000  0.0000000000000000
  0.5000000000000000  0.5000000000000000  0.0000000000000000
  0.8272803098748985  0.4538044740481856  0.5000000000000000
  0.1727196901251014  0.5461955259518143  0.5000000000000000
  0.6727196901251015  0.9538044740481857  0.5000000000000000
  0.3272803098748985  0.0461955259518144  0.5000000000000000
 
  0.00000000E+00  0.00000000E+00  0.00000000E+00
  0.00000000E+00  0.00000000E+00  0.00000000E+00
  0.00000000E+00  0.00000000E+00  0.00000000E+00
  0.00000000E+00  0.00000000E+00  0.00000000E+00
  0.00000000E+00  0.00000000E+00  0.00000000E+00
  0.00000000E+00  0.00000000E+00  0.00000000E+00
  0.00000000E+00  0.00000000E+00  0.00000000E+00
  0.00000000E+00  0.00000000E+00  0.00000000E+00
  0.00000000E+00  0.00000000E+00  0.00000000E+00
  0.00000000E+00  0.00000000E+00  0.00000000E+00
  0.00000000E+00  0.00000000E+00  0.00000000E+00
  0.00000000E+00  0.00000000E+00  0.00000000E+00
  0.00000000E+00  0.00000000E+00  0.00000000E+00
  0.00000000E+00  0.00000000E+00  0.00000000E+00
\end{verbatim}
\end{footnotesize}

\subsubsection*{Pbam -- 50~GPa}
\begin{footnotesize}
\begin{verbatim}
Al4C3 Pbam polymorph (findsym-output) - 
   1.00000000000000     
     8.4823493958000000    0.0000000000000000    0.0000000000000000
     0.0000000000000000    4.8977818489000002    0.0000000000000000
     0.0000000000000000    0.0000000000000000    2.9013054371000000
   Al   C 
     8     6
Direct
  0.7858800289999976  0.1601099969999993  0.0000000000000000
  0.2141199710000024  0.8398900030000007  0.0000000000000000
  0.7141199710000024  0.6601099969999993  0.0000000000000000
  0.2858800289999976  0.3398900030000007  0.0000000000000000
  0.0310999989999985  0.2530300020000027  0.5000000000000000
  0.9689000250000035  0.7469699979999973  0.5000000000000000
  0.4688999950000010  0.7530300020000027  0.5000000000000000
  0.5310999749999965  0.2469699979999973  0.5000000000000000
  0.0000000000000000  0.0000000000000000  0.0000000000000000
  0.5000000000000000  0.5000000000000000  0.0000000000000000
  0.8245900269999993  0.4467200039999994  0.5000000000000000
  0.1754099730000007  0.5532799960000006  0.5000000000000000
  0.6754099730000007  0.9467200039999994  0.5000000000000000
  0.3245900269999993  0.0532799960000006  0.5000000000000000
 
  0.00000000E+00  0.00000000E+00  0.00000000E+00
  0.00000000E+00  0.00000000E+00  0.00000000E+00
  0.00000000E+00  0.00000000E+00  0.00000000E+00
  0.00000000E+00  0.00000000E+00  0.00000000E+00
  0.00000000E+00  0.00000000E+00  0.00000000E+00
  0.00000000E+00  0.00000000E+00  0.00000000E+00
  0.00000000E+00  0.00000000E+00  0.00000000E+00
  0.00000000E+00  0.00000000E+00  0.00000000E+00
  0.00000000E+00  0.00000000E+00  0.00000000E+00
  0.00000000E+00  0.00000000E+00  0.00000000E+00
  0.00000000E+00  0.00000000E+00  0.00000000E+00
  0.00000000E+00  0.00000000E+00  0.00000000E+00
  0.00000000E+00  0.00000000E+00  0.00000000E+00
  0.00000000E+00  0.00000000E+00  0.00000000E+00
\end{verbatim}
\end{footnotesize}

\subsubsection*{Pnma -- 0~GPa}
\begin{footnotesize}
\begin{verbatim}
Al4C3 Pnma (POSCAR.IV, symmetrized via s
   1.00000000000000     
     8.8687877986054335    0.0000000000000000    0.0000000000000000
     0.0000000000000000    3.0595706192262035    0.0000000000000000
     0.0000000000000000    0.0000000000000000   10.5424912694858808
   Al   C 
    16    12
Direct
  0.2735502672922275  0.2500000000000000  0.3075329915271308
  0.2264497327077724  0.7500000000000000  0.8075329915271308
  0.7264497327077724  0.7500000000000000  0.6924670084728692
  0.7735502672922276  0.2500000000000000  0.1924670084728693
  0.3615341662279118  0.2500000000000000  0.0187587654397753
  0.1384658337720882  0.7500000000000000  0.5187587654397752
  0.6384658337720881  0.7500000000000000  0.9812412345602248
  0.8615341662279119  0.2500000000000000  0.4812412345602247
  0.5160590391368277  0.2500000000000000  0.7948720507037115
  0.9839409608631723  0.7500000000000000  0.2948720507037115
  0.4839409608631723  0.7500000000000000  0.2051279492962884
  0.0160590391368277  0.2500000000000000  0.7051279492962885
  0.0813324512462430  0.2500000000000000  0.0833467341491087
  0.4186675487537640  0.7500000000000000  0.5833467341491085
  0.9186675487537570  0.7500000000000000  0.9166532658508915
  0.5813324512462430  0.2500000000000000  0.4166532658508914
  0.2457612108131275  0.2500000000000000  0.6385376537495265
  0.2542387891868724  0.7500000000000000  0.1385376537495265
  0.7542387891868727  0.7500000000000000  0.3614623462504735
  0.7457612108131273  0.2500000000000000  0.8614623462504735
  0.5875004958432158  0.2500000000000000  0.6136101721742773
  0.9124995041567842  0.7500000000000000  0.1136101721742773
  0.4124995041567915  0.7500000000000000  0.3863898278257227
  0.0875004958432156  0.2500000000000000  0.8863898278257227
  0.0722337674812502  0.2500000000000000  0.3872433425577897
  0.4277662325187498  0.7500000000000000  0.8872433425577972
  0.9277662325187502  0.7500000000000000  0.6127566574422028
  0.5722337674812498  0.2500000000000000  0.1127566574422030
 
  0.00000000E+00  0.00000000E+00  0.00000000E+00
  0.00000000E+00  0.00000000E+00  0.00000000E+00
  0.00000000E+00  0.00000000E+00  0.00000000E+00
  0.00000000E+00  0.00000000E+00  0.00000000E+00
  0.00000000E+00  0.00000000E+00  0.00000000E+00
  0.00000000E+00  0.00000000E+00  0.00000000E+00
  0.00000000E+00  0.00000000E+00  0.00000000E+00
  0.00000000E+00  0.00000000E+00  0.00000000E+00
  0.00000000E+00  0.00000000E+00  0.00000000E+00
  0.00000000E+00  0.00000000E+00  0.00000000E+00
  0.00000000E+00  0.00000000E+00  0.00000000E+00
  0.00000000E+00  0.00000000E+00  0.00000000E+00
  0.00000000E+00  0.00000000E+00  0.00000000E+00
  0.00000000E+00  0.00000000E+00  0.00000000E+00
  0.00000000E+00  0.00000000E+00  0.00000000E+00
  0.00000000E+00  0.00000000E+00  0.00000000E+00
  0.00000000E+00  0.00000000E+00  0.00000000E+00
  0.00000000E+00  0.00000000E+00  0.00000000E+00
  0.00000000E+00  0.00000000E+00  0.00000000E+00
  0.00000000E+00  0.00000000E+00  0.00000000E+00
  0.00000000E+00  0.00000000E+00  0.00000000E+00
  0.00000000E+00  0.00000000E+00  0.00000000E+00
  0.00000000E+00  0.00000000E+00  0.00000000E+00
  0.00000000E+00  0.00000000E+00  0.00000000E+00
  0.00000000E+00  0.00000000E+00  0.00000000E+00
  0.00000000E+00  0.00000000E+00  0.00000000E+00
  0.00000000E+00  0.00000000E+00  0.00000000E+00
  0.00000000E+00  0.00000000E+00  0.00000000E+00
\end{verbatim}
\end{footnotesize}

\subsubsection*{Pnma -- 50~GPa}
\begin{footnotesize}
\begin{verbatim}
Al4C3 Pnma (POSCAR.IV, symmetrized via s
   1.00000000000000     
     8.2890294394235049    0.0000000000000000    0.0000000000000000
     0.0000000000000000    2.8595645211794394    0.0000000000000000
     0.0000000000000000    0.0000000000000000    9.8533218385690624
   Al   C 
    16    12
Direct
  0.2629312655562072  0.2500000000000000  0.2981934154944613
  0.2370687344437928  0.7500000000000000  0.7981934154944613
  0.7370687344437928  0.7500000000000000  0.7018065845055387
  0.7629312655562072  0.2500000000000000  0.2018065845055387
  0.3469030433064250  0.2500000000000000  0.0112167087473821
  0.1530969566935750  0.7500000000000000  0.5112167087473821
  0.6530969566935749  0.7500000000000000  0.9887832912526179
  0.8469030433064251  0.2500000000000000  0.4887832912526179
  0.5225009079149764  0.2500000000000000  0.7942536374118344
  0.9774990920850236  0.7500000000000000  0.2942536374118344
  0.4774990920850237  0.7500000000000000  0.2057463625881656
  0.0225009079149763  0.2500000000000000  0.7057463625881656
  0.0689590485967116  0.2500000000000000  0.0894805905942223
  0.4310409514032955  0.7500000000000000  0.5894805905942223
  0.9310409514032884  0.7500000000000000  0.9105194094057777
  0.5689590485967116  0.2500000000000000  0.4105194094057777
  0.2592235199889221  0.2500000000000000  0.6419226335859984
  0.2407764800110779  0.7500000000000000  0.1419226335859984
  0.7407764800110779  0.7500000000000000  0.3580773664140016
  0.7592235199889221  0.2500000000000000  0.8580773664140016
  0.6043531701401923  0.2500000000000000  0.6081829935710156
  0.8956468298598077  0.7500000000000000  0.1081829935710156
  0.3956468298598077  0.7500000000000000  0.3918170064289844
  0.1043531701401923  0.2500000000000000  0.8918170064289844
  0.0641760130956121  0.2500000000000000  0.3891272411387470
  0.4358239869043878  0.7500000000000000  0.8891272411387470
  0.9358239869043878  0.7500000000000000  0.6108727588612530
  0.5641760130956122  0.2500000000000000  0.1108727588612530
 
  0.00000000E+00  0.00000000E+00  0.00000000E+00
  0.00000000E+00  0.00000000E+00  0.00000000E+00
  0.00000000E+00  0.00000000E+00  0.00000000E+00
  0.00000000E+00  0.00000000E+00  0.00000000E+00
  0.00000000E+00  0.00000000E+00  0.00000000E+00
  0.00000000E+00  0.00000000E+00  0.00000000E+00
  0.00000000E+00  0.00000000E+00  0.00000000E+00
  0.00000000E+00  0.00000000E+00  0.00000000E+00
  0.00000000E+00  0.00000000E+00  0.00000000E+00
  0.00000000E+00  0.00000000E+00  0.00000000E+00
  0.00000000E+00  0.00000000E+00  0.00000000E+00
  0.00000000E+00  0.00000000E+00  0.00000000E+00
  0.00000000E+00  0.00000000E+00  0.00000000E+00
  0.00000000E+00  0.00000000E+00  0.00000000E+00
  0.00000000E+00  0.00000000E+00  0.00000000E+00
  0.00000000E+00  0.00000000E+00  0.00000000E+00
  0.00000000E+00  0.00000000E+00  0.00000000E+00
  0.00000000E+00  0.00000000E+00  0.00000000E+00
  0.00000000E+00  0.00000000E+00  0.00000000E+00
  0.00000000E+00  0.00000000E+00  0.00000000E+00
  0.00000000E+00  0.00000000E+00  0.00000000E+00
  0.00000000E+00  0.00000000E+00  0.00000000E+00
  0.00000000E+00  0.00000000E+00  0.00000000E+00
  0.00000000E+00  0.00000000E+00  0.00000000E+00
  0.00000000E+00  0.00000000E+00  0.00000000E+00
  0.00000000E+00  0.00000000E+00  0.00000000E+00
  0.00000000E+00  0.00000000E+00  0.00000000E+00
  0.00000000E+00  0.00000000E+00  0.00000000E+00
\end{verbatim}
\end{footnotesize}

\subsubsection*{R-3m -- 0~GPa}
\begin{footnotesize}
\begin{verbatim}
Al4C3 R-3m polymorph (findsym-output, fr
   1.00000000000000     
     3.3554638352239903    0.0000000000000000    0.0000000000000000
    -1.6777319176119951    2.9059169228212531    0.0000000000000000
     0.0000000000000000    0.0000000000000000   25.0854610045852198
   Al   C 
    12     9
Direct
  0.0000000000000000 -0.0000000000000000  0.6298944952465831
 -0.0000000000000000  0.0000000000000000  0.3701055047534167
  0.6666666870000029  0.3333333429999996  0.9632278082465803
  0.6666666870000029  0.3333333429999996  0.7034388187534176
  0.3333333429999996  0.6666666870000029  0.2965611512465801
  0.3333333429999996  0.6666666870000029  0.0367721727534216
  0.0000000000000000 -0.0000000000000000  0.2065187870111717
 -0.0000000000000000  0.0000000000000000  0.7934812129888283
  0.6666666870000029  0.3333333429999996  0.5398521010111724
  0.6666666870000029  0.3333333429999996  0.1268145409888303
  0.3333333429999996  0.6666666870000029  0.8731854740111745
  0.3333333429999996  0.6666666870000029  0.4601478689888250
  0.0000000000000000  0.0000000000000000  0.7167106476763591
  0.0000000000000000 -0.0000000000000000  0.2832893523236410
  0.6666666870000029  0.3333333429999996  0.0500439796763613
  0.6666666870000029  0.3333333429999996  0.6166226653236381
  0.3333333429999996  0.6666666870000029  0.3833773046763594
  0.3333333429999996  0.6666666870000029  0.9499560393236367
  0.0000000000000000 -0.0000000000000000  0.5000000000000000
  0.6666666870000029  0.3333333429999996  0.8333333129999971
  0.3333333429999996  0.6666666870000029  0.1666666719999981
 
  0.00000000E+00  0.00000000E+00  0.00000000E+00
  0.00000000E+00  0.00000000E+00  0.00000000E+00
  0.00000000E+00  0.00000000E+00  0.00000000E+00
  0.00000000E+00  0.00000000E+00  0.00000000E+00
  0.00000000E+00  0.00000000E+00  0.00000000E+00
  0.00000000E+00  0.00000000E+00  0.00000000E+00
  0.00000000E+00  0.00000000E+00  0.00000000E+00
  0.00000000E+00  0.00000000E+00  0.00000000E+00
  0.00000000E+00  0.00000000E+00  0.00000000E+00
  0.00000000E+00  0.00000000E+00  0.00000000E+00
  0.00000000E+00  0.00000000E+00  0.00000000E+00
  0.00000000E+00  0.00000000E+00  0.00000000E+00
  0.00000000E+00  0.00000000E+00  0.00000000E+00
  0.00000000E+00  0.00000000E+00  0.00000000E+00
  0.00000000E+00  0.00000000E+00  0.00000000E+00
  0.00000000E+00  0.00000000E+00  0.00000000E+00
  0.00000000E+00  0.00000000E+00  0.00000000E+00
  0.00000000E+00  0.00000000E+00  0.00000000E+00
  0.00000000E+00  0.00000000E+00  0.00000000E+00
  0.00000000E+00  0.00000000E+00  0.00000000E+00
  0.00000000E+00  0.00000000E+00  0.00000000E+00
\end{verbatim}
\end{footnotesize}

\subsubsection*{R-3m -- 50~GPa}
\begin{footnotesize}
\begin{verbatim}
Al4C3 R-3m polymorph (findsym-output, fr
   1.00000000000000     
     3.1293263730393432    0.0000000000000000    0.0000000000000000
    -1.5646631865196716    2.7100761358194903    0.0000000000000000
     0.0000000000000000    0.0000000000000000   23.3948564360725193
   Al   C 
    12     9
Direct
  0.0000000000000000  0.0000000000000000  0.6295050301620364
  0.0000000000000000  0.0000000000000000  0.3704949698379636
  0.6666666870000029  0.3333333429999996  0.9628383431620335
  0.6666666870000029  0.3333333429999996  0.7038282838379644
  0.3333333429999996  0.6666666870000029  0.2961716861620332
  0.3333333429999996  0.6666666870000029  0.0371616378379684
  0.0000000000000000  0.0000000000000000  0.2060826556182036
  0.0000000000000000  0.0000000000000000  0.7939173443817964
  0.6666666870000029  0.3333333429999996  0.5394159696182044
  0.6666666870000029  0.3333333429999996  0.1272506723817983
  0.3333333429999996  0.6666666870000029  0.8727493426182065
  0.3333333429999996  0.6666666870000029  0.4605840003817931
  0.0000000000000000  0.0000000000000000  0.7166862778946808
  0.0000000000000000  0.0000000000000000  0.2833137221053192
  0.6666666870000029  0.3333333429999996  0.0500196098946830
  0.6666666870000029  0.3333333429999996  0.6166470351053164
  0.3333333429999996  0.6666666870000029  0.3833529348946811
  0.3333333429999996  0.6666666870000029  0.9499804091053150
  0.0000000000000000  0.0000000000000000  0.5000000000000000
  0.6666666870000029  0.3333333429999996  0.8333333129999971
  0.3333333429999996  0.6666666870000029  0.1666666719999981
 
  0.00000000E+00  0.00000000E+00  0.00000000E+00
  0.00000000E+00  0.00000000E+00  0.00000000E+00
  0.00000000E+00  0.00000000E+00  0.00000000E+00
  0.00000000E+00  0.00000000E+00  0.00000000E+00
  0.00000000E+00  0.00000000E+00  0.00000000E+00
  0.00000000E+00  0.00000000E+00  0.00000000E+00
  0.00000000E+00  0.00000000E+00  0.00000000E+00
  0.00000000E+00  0.00000000E+00  0.00000000E+00
  0.00000000E+00  0.00000000E+00  0.00000000E+00
  0.00000000E+00  0.00000000E+00  0.00000000E+00
  0.00000000E+00  0.00000000E+00  0.00000000E+00
  0.00000000E+00  0.00000000E+00  0.00000000E+00
  0.00000000E+00  0.00000000E+00  0.00000000E+00
  0.00000000E+00  0.00000000E+00  0.00000000E+00
  0.00000000E+00  0.00000000E+00  0.00000000E+00
  0.00000000E+00  0.00000000E+00  0.00000000E+00
  0.00000000E+00  0.00000000E+00  0.00000000E+00
  0.00000000E+00  0.00000000E+00  0.00000000E+00
  0.00000000E+00  0.00000000E+00  0.00000000E+00
  0.00000000E+00  0.00000000E+00  0.00000000E+00
  0.00000000E+00  0.00000000E+00  0.00000000E+00
\end{verbatim}
\end{footnotesize}

\subsubsection*{Pnma-III -- 0~GPa}
\begin{footnotesize}
\begin{verbatim}
Al4C3 Pnma (POSCAR.III, symmetrized via 
   1.00000000000000     
     9.3185950946021627    0.0000000000000000    0.0000000000000000
     0.0000000000000000    3.1424134406869162    0.0000000000000000
     0.0000000000000000    0.0000000000000000   10.1601230062129950
   Al   C 
    16    12
Direct
  0.7449744616347906  0.2500000000000000  0.7201669400270182
  0.7550255383652094  0.7500000000000000  0.2201669400270180
  0.2550255383652096  0.7500000000000000  0.2798330599729819
  0.2449744616347904  0.2500000000000000  0.7798330599729818
  0.6082545802170108  0.2500000000000000  0.4786801871503972
  0.8917454197829892  0.7500000000000000  0.9786801871503900
  0.3917454197829891  0.7500000000000000  0.5213198128496100
  0.1082545802170108  0.2500000000000000  0.0213198128496100
  0.1325614982528854  0.2500000000000000  0.5176342081888706
  0.3674385017471147  0.7500000000000000  0.0176342081888705
  0.8674385017471147  0.7500000000000000  0.4823657918111295
  0.6325614982528853  0.2500000000000000  0.9823657918111294
  0.9830998624780699  0.2500000000000000  0.2823658008199360
  0.5169001375219301  0.7500000000000000  0.7823658008199362
  0.0169001375219303  0.7500000000000000  0.7176341991800638
  0.4830998624780697  0.2500000000000000  0.2176341991800640
  0.9307094114206790  0.2500000000000000  0.6227096332100521
  0.5692905885793210  0.7500000000000000  0.1227096332100522
  0.0692905885793140  0.7500000000000000  0.3772903667899479
  0.4307094114206860  0.2500000000000000  0.8772903667899479
  0.3977891753639959  0.2500000000000000  0.3993912080842171
  0.1022108246360040  0.7500000000000000  0.8993912080842100
  0.6022108246360039  0.7500000000000000  0.6006087919157900
  0.8977891753639961  0.2500000000000000  0.1006087919157901
  0.7741899310421310  0.2500000000000000  0.3593595510598142
  0.7258100689578690  0.7500000000000000  0.8593595510598141
  0.2258100689578688  0.7500000000000000  0.6406404489401859
  0.2741899310421312  0.2500000000000000  0.1406404489401859
 
  0.00000000E+00  0.00000000E+00  0.00000000E+00
  0.00000000E+00  0.00000000E+00  0.00000000E+00
  0.00000000E+00  0.00000000E+00  0.00000000E+00
  0.00000000E+00  0.00000000E+00  0.00000000E+00
  0.00000000E+00  0.00000000E+00  0.00000000E+00
  0.00000000E+00  0.00000000E+00  0.00000000E+00
  0.00000000E+00  0.00000000E+00  0.00000000E+00
  0.00000000E+00  0.00000000E+00  0.00000000E+00
  0.00000000E+00  0.00000000E+00  0.00000000E+00
  0.00000000E+00  0.00000000E+00  0.00000000E+00
  0.00000000E+00  0.00000000E+00  0.00000000E+00
  0.00000000E+00  0.00000000E+00  0.00000000E+00
  0.00000000E+00  0.00000000E+00  0.00000000E+00
  0.00000000E+00  0.00000000E+00  0.00000000E+00
  0.00000000E+00  0.00000000E+00  0.00000000E+00
  0.00000000E+00  0.00000000E+00  0.00000000E+00
  0.00000000E+00  0.00000000E+00  0.00000000E+00
  0.00000000E+00  0.00000000E+00  0.00000000E+00
  0.00000000E+00  0.00000000E+00  0.00000000E+00
  0.00000000E+00  0.00000000E+00  0.00000000E+00
  0.00000000E+00  0.00000000E+00  0.00000000E+00
  0.00000000E+00  0.00000000E+00  0.00000000E+00
  0.00000000E+00  0.00000000E+00  0.00000000E+00
  0.00000000E+00  0.00000000E+00  0.00000000E+00
  0.00000000E+00  0.00000000E+00  0.00000000E+00
  0.00000000E+00  0.00000000E+00  0.00000000E+00
  0.00000000E+00  0.00000000E+00  0.00000000E+00
  0.00000000E+00  0.00000000E+00  0.00000000E+00
\end{verbatim}
\end{footnotesize}

\subsubsection*{Pnma-III -- 50~GPa}
\begin{footnotesize}
\begin{verbatim}
Al4C3 Pnma (POSCAR.III, symmetrized via 
   1.00000000000000     
     8.7032434407543473    0.0000000000000000    0.0000000000000000
     0.0000000000000000    2.9349047670972270    0.0000000000000000
     0.0000000000000000    0.0000000000000000    9.4892012168552959
   Al   C 
    16    12
Direct
  0.7372093599188929  0.2500000000000000  0.7136306477819190
  0.7627906400811071  0.7500000000000000  0.2136306477819190
  0.2627906400811071  0.7500000000000000  0.2863693522180810
  0.2372093599188929  0.2500000000000000  0.7863693522180810
  0.6063043166502808  0.2500000000000000  0.4707003951628863
  0.8936956833497192  0.7500000000000000  0.9707003951628863
  0.3936956833497192  0.7500000000000000  0.5292996048371137
  0.1063043166502808  0.2500000000000000  0.0292996048371137
  0.1374957759450339  0.2500000000000000  0.5177890409486778
  0.3625042240549661  0.7500000000000000  0.0177890409486778
  0.8625042240549661  0.7500000000000000  0.4822109590513222
  0.6374957759450339  0.2500000000000000  0.9822109590513222
  0.9932809183437712  0.2500000000000000  0.2795179626542890
  0.5067190816562288  0.7500000000000000  0.7795179626542890
  0.0067190816562288  0.7500000000000000  0.7204820373457110
  0.4932809183437712  0.2500000000000000  0.2204820373457110
  0.9285470810138605  0.2500000000000000  0.6187220005827854
  0.5714529189861395  0.7500000000000000  0.1187220005827854
  0.0714529189861395  0.7500000000000000  0.3812779994172146
  0.4285470810138605  0.2500000000000000  0.8812779994172146
  0.3911565475003371  0.2500000000000000  0.4018945523721129
  0.1088434524996629  0.7500000000000000  0.9018945523721058
  0.6088434524996629  0.7500000000000000  0.5981054476278942
  0.8911565475003371  0.2500000000000000  0.0981054476278942
  0.7744048364141989  0.2500000000000000  0.3539263382519593
  0.7255951635858011  0.7500000000000000  0.8539263382519593
  0.2255951635858011  0.7500000000000000  0.6460736617480407
  0.2744048364141989  0.2500000000000000  0.1460736617480407
 
  0.00000000E+00  0.00000000E+00  0.00000000E+00
  0.00000000E+00  0.00000000E+00  0.00000000E+00
  0.00000000E+00  0.00000000E+00  0.00000000E+00
  0.00000000E+00  0.00000000E+00  0.00000000E+00
  0.00000000E+00  0.00000000E+00  0.00000000E+00
  0.00000000E+00  0.00000000E+00  0.00000000E+00
  0.00000000E+00  0.00000000E+00  0.00000000E+00
  0.00000000E+00  0.00000000E+00  0.00000000E+00
  0.00000000E+00  0.00000000E+00  0.00000000E+00
  0.00000000E+00  0.00000000E+00  0.00000000E+00
  0.00000000E+00  0.00000000E+00  0.00000000E+00
  0.00000000E+00  0.00000000E+00  0.00000000E+00
  0.00000000E+00  0.00000000E+00  0.00000000E+00
  0.00000000E+00  0.00000000E+00  0.00000000E+00
  0.00000000E+00  0.00000000E+00  0.00000000E+00
  0.00000000E+00  0.00000000E+00  0.00000000E+00
  0.00000000E+00  0.00000000E+00  0.00000000E+00
  0.00000000E+00  0.00000000E+00  0.00000000E+00
  0.00000000E+00  0.00000000E+00  0.00000000E+00
  0.00000000E+00  0.00000000E+00  0.00000000E+00
  0.00000000E+00  0.00000000E+00  0.00000000E+00
  0.00000000E+00  0.00000000E+00  0.00000000E+00
  0.00000000E+00  0.00000000E+00  0.00000000E+00
  0.00000000E+00  0.00000000E+00  0.00000000E+00
  0.00000000E+00  0.00000000E+00  0.00000000E+00
  0.00000000E+00  0.00000000E+00  0.00000000E+00
  0.00000000E+00  0.00000000E+00  0.00000000E+00
  0.00000000E+00  0.00000000E+00  0.00000000E+00
\end{verbatim}
\end{footnotesize}

\subsubsection*{anti-CaFe$_2$O$_4$ -- 0~GPa}
\begin{footnotesize}
\begin{verbatim}
Al4C3 anti-CaFe2O4-type, volume pre-scal
   1.00000000000000     
     2.9852194567723527    0.0000000000000000    0.0000000000000000
     0.0000000000000000    9.1198705024020921    0.0000000000000000
     0.0000000000000000    0.0000000000000000   10.5792008400718061
   Al   C 
    16    12
Direct
  0.2500000000000000  0.4144072947002274  0.5809008803943491
  0.7500000000000000  0.5855927052997728  0.4190991196056512
  0.2500000000000000  0.9144072947002272  0.9190991196056509
  0.7500000000000000  0.0855927052997727  0.0809008803943488
  0.2500000000000000  0.4889137898529304  0.2054812901082441
  0.7500000000000000  0.5110862101470698  0.7945187098917560
  0.2500000000000000  0.9889137898529302  0.2945187098917558
  0.7500000000000000  0.0110862101470698  0.7054812901082440
  0.2500000000000000  0.1351843018774951  0.5212377979060971
  0.7500000000000000  0.8648156981224979  0.4787622020939027
  0.2500000000000000  0.6351843018775021  0.9787622020939029
  0.7500000000000000  0.3648156981225048  0.0212377979060973
  0.7500000000000000  0.2770966780917617  0.3071244884571125
  0.2500000000000000  0.7229033219082382  0.6928755115428877
  0.2500000000000000  0.2229033219082383  0.8071244884571123
  0.7500000000000000  0.7770966780917618  0.1928755115428875
  0.7500000000000000  0.7411535668307502  0.8617177029566370
  0.2500000000000000  0.2588464331692499  0.1382822970433632
  0.7500000000000000  0.2411535668307501  0.6382822970433630
  0.2500000000000000  0.7588464331692498  0.3617177029566369
  0.2500000000000000  0.9221420520041845  0.6119178216534967
  0.7500000000000000  0.5778579479958155  0.1119178216534966
  0.2500000000000000  0.4221420520041844  0.8880821783465033
  0.7500000000000000  0.0778579479958156  0.3880821783465034
  0.7500000000000000  0.0805052220197591  0.8841128946722824
  0.2500000000000000  0.9194947779802410  0.1158871053277176
  0.7500000000000000  0.5805052220197590  0.6158871053277176
  0.2500000000000000  0.4194947779802408  0.3841128946722825
 
  0.00000000E+00  0.00000000E+00  0.00000000E+00
  0.00000000E+00  0.00000000E+00  0.00000000E+00
  0.00000000E+00  0.00000000E+00  0.00000000E+00
  0.00000000E+00  0.00000000E+00  0.00000000E+00
  0.00000000E+00  0.00000000E+00  0.00000000E+00
  0.00000000E+00  0.00000000E+00  0.00000000E+00
  0.00000000E+00  0.00000000E+00  0.00000000E+00
  0.00000000E+00  0.00000000E+00  0.00000000E+00
  0.00000000E+00  0.00000000E+00  0.00000000E+00
  0.00000000E+00  0.00000000E+00  0.00000000E+00
  0.00000000E+00  0.00000000E+00  0.00000000E+00
  0.00000000E+00  0.00000000E+00  0.00000000E+00
  0.00000000E+00  0.00000000E+00  0.00000000E+00
  0.00000000E+00  0.00000000E+00  0.00000000E+00
  0.00000000E+00  0.00000000E+00  0.00000000E+00
  0.00000000E+00  0.00000000E+00  0.00000000E+00
  0.00000000E+00  0.00000000E+00  0.00000000E+00
  0.00000000E+00  0.00000000E+00  0.00000000E+00
  0.00000000E+00  0.00000000E+00  0.00000000E+00
  0.00000000E+00  0.00000000E+00  0.00000000E+00
  0.00000000E+00  0.00000000E+00  0.00000000E+00
  0.00000000E+00  0.00000000E+00  0.00000000E+00
  0.00000000E+00  0.00000000E+00  0.00000000E+00
  0.00000000E+00  0.00000000E+00  0.00000000E+00
  0.00000000E+00  0.00000000E+00  0.00000000E+00
  0.00000000E+00  0.00000000E+00  0.00000000E+00
  0.00000000E+00  0.00000000E+00  0.00000000E+00
  0.00000000E+00  0.00000000E+00  0.00000000E+00
\end{verbatim}
\end{footnotesize}

\subsubsection*{anti-CaFe$_2$O$_4$ -- 50~GPa}
\begin{footnotesize}
\begin{verbatim}
Al4C3 anti-CaFe2O4-type, volume pre-scal
   1.00000000000000     
     2.7900451726561628    0.0000000000000000    0.0000000000000000
     0.0000000000000000    8.5236114258706834    0.0000000000000000
     0.0000000000000000    0.0000000000000000    9.8875304351378670
   Al   C 
    16    12
Direct
  0.2500000000000000  0.4245267340863990  0.5873716912218541
  0.7500000000000000  0.5754732659136010  0.4126283087781459
  0.2500000000000000  0.9245267340863990  0.9126283087781459
  0.7500000000000000  0.0754732659136010  0.0873716912218541
  0.2500000000000000  0.4800601400098898  0.2061656728190329
  0.7500000000000000  0.5199398599901102  0.7938343271809671
  0.2500000000000000  0.9800601400098898  0.2938343271809671
  0.7500000000000000  0.0199398599901102  0.7061656728190329
  0.2500000000000000  0.1484441334435687  0.5143475653135141
  0.7500000000000000  0.8515558665564242  0.4856524346864859
  0.2500000000000000  0.6484441334435758  0.9856524346864859
  0.7500000000000000  0.3515558665564313  0.0143475653135141
  0.7500000000000000  0.2662336654069293  0.3000390016478320
  0.2500000000000000  0.7337663345930707  0.6999609983521680
  0.2500000000000000  0.2337663345930707  0.8000390016478320
  0.7500000000000000  0.7662336654069293  0.1999609983521680
  0.7500000000000000  0.7536300515530598  0.8590001755942538
  0.2500000000000000  0.2463699484469402  0.1409998244057462
  0.7500000000000000  0.2536300515530598  0.6409998244057462
  0.2500000000000000  0.7463699484469402  0.3590001755942538
  0.2500000000000000  0.9301796991213749  0.6107702309581455
  0.7500000000000000  0.5698203008786251  0.1107702309581455
  0.2500000000000000  0.4301796991213749  0.8892297690418545
  0.7500000000000000  0.0698203008786251  0.3892297690418545
  0.7500000000000000  0.0968949186839723  0.8893454137177415
  0.2500000000000000  0.9031050813160277  0.1106545862822585
  0.7500000000000000  0.5968949186839723  0.6106545862822585
  0.2500000000000000  0.4031050813160277  0.3893454137177415
 
  0.00000000E+00  0.00000000E+00  0.00000000E+00
  0.00000000E+00  0.00000000E+00  0.00000000E+00
  0.00000000E+00  0.00000000E+00  0.00000000E+00
  0.00000000E+00  0.00000000E+00  0.00000000E+00
  0.00000000E+00  0.00000000E+00  0.00000000E+00
  0.00000000E+00  0.00000000E+00  0.00000000E+00
  0.00000000E+00  0.00000000E+00  0.00000000E+00
  0.00000000E+00  0.00000000E+00  0.00000000E+00
  0.00000000E+00  0.00000000E+00  0.00000000E+00
  0.00000000E+00  0.00000000E+00  0.00000000E+00
  0.00000000E+00  0.00000000E+00  0.00000000E+00
  0.00000000E+00  0.00000000E+00  0.00000000E+00
  0.00000000E+00  0.00000000E+00  0.00000000E+00
  0.00000000E+00  0.00000000E+00  0.00000000E+00
  0.00000000E+00  0.00000000E+00  0.00000000E+00
  0.00000000E+00  0.00000000E+00  0.00000000E+00
  0.00000000E+00  0.00000000E+00  0.00000000E+00
  0.00000000E+00  0.00000000E+00  0.00000000E+00
  0.00000000E+00  0.00000000E+00  0.00000000E+00
  0.00000000E+00  0.00000000E+00  0.00000000E+00
  0.00000000E+00  0.00000000E+00  0.00000000E+00
  0.00000000E+00  0.00000000E+00  0.00000000E+00
  0.00000000E+00  0.00000000E+00  0.00000000E+00
  0.00000000E+00  0.00000000E+00  0.00000000E+00
  0.00000000E+00  0.00000000E+00  0.00000000E+00
  0.00000000E+00  0.00000000E+00  0.00000000E+00
  0.00000000E+00  0.00000000E+00  0.00000000E+00
  0.00000000E+00  0.00000000E+00  0.00000000E+00
\end{verbatim}
\end{footnotesize}

\subsubsection*{I-43d -- 0~GPa}
\begin{footnotesize}
\begin{verbatim}
Al4C3 I-43d polymorph (findsym-output, f
   1.00000000000000     
     6.5463773857702332    0.0000000000000000    0.0000000000000000
     0.0000000000000000    6.5463773857702332    0.0000000000000000
     0.0000000000000000    0.0000000000000000    6.5463773857702332
   Al   C 
    16    12
Direct
  0.2990010309058072  0.2990010309058072  0.2990010309058072
  0.2009989690941928  0.7009989690941930  0.7990010309058070
  0.7009989690941930  0.7990010309058070  0.2009989690941928
  0.7990010309058070  0.2009989690941928  0.7009989690941930
  0.5490010309058070  0.5490010309058070  0.5490010309058070
  0.9509989690941930  0.4509989690941928  0.0490010309058072
  0.0490010309058072  0.9509989690941930  0.4509989690941928
  0.4509989690941928  0.0490010309058072  0.9509989690941930
  0.7990010309058070  0.7990010309058070  0.7990010309058070
  0.7009989690941930  0.2009989690941928  0.2990010309058072
  0.2009989690941928  0.2990010309058072  0.7009989690941930
  0.2990010309058072  0.7009989690941930  0.2009989690941928
  0.0490010309058072  0.0490010309058072  0.0490010309058072
  0.4509989690941928  0.9509989690941930  0.5490010309058070
  0.5490010309058070  0.4509989690941928  0.9509989690941930
  0.9509989690941930  0.5490010309058070  0.4509989690941928
  0.3750000000000000 -0.0000000000000000  0.2500000000000000
  0.1250000000000000 -0.0000000000000000  0.7500000000000000
  0.6250000000000000  0.5000000000000000  0.2500000000000000
  0.8750000000000000  0.5000000000000000  0.7500000000000000
  0.2500000000000000  0.3750000000000000 -0.0000000000000000
  0.7500000000000000  0.1250000000000000 -0.0000000000000000
  0.2500000000000000  0.6250000000000000  0.5000000000000000
  0.7500000000000000  0.8750000000000000  0.5000000000000000
 -0.0000000000000000  0.2500000000000000  0.3750000000000000
 -0.0000000000000000  0.7500000000000000  0.1250000000000000
  0.5000000000000000  0.2500000000000000  0.6250000000000000
  0.5000000000000000  0.7500000000000000  0.8750000000000000
 
  0.00000000E+00  0.00000000E+00  0.00000000E+00
  0.00000000E+00  0.00000000E+00  0.00000000E+00
  0.00000000E+00  0.00000000E+00  0.00000000E+00
  0.00000000E+00  0.00000000E+00  0.00000000E+00
  0.00000000E+00  0.00000000E+00  0.00000000E+00
  0.00000000E+00  0.00000000E+00  0.00000000E+00
  0.00000000E+00  0.00000000E+00  0.00000000E+00
  0.00000000E+00  0.00000000E+00  0.00000000E+00
  0.00000000E+00  0.00000000E+00  0.00000000E+00
  0.00000000E+00  0.00000000E+00  0.00000000E+00
  0.00000000E+00  0.00000000E+00  0.00000000E+00
  0.00000000E+00  0.00000000E+00  0.00000000E+00
  0.00000000E+00  0.00000000E+00  0.00000000E+00
  0.00000000E+00  0.00000000E+00  0.00000000E+00
  0.00000000E+00  0.00000000E+00  0.00000000E+00
  0.00000000E+00  0.00000000E+00  0.00000000E+00
  0.00000000E+00  0.00000000E+00  0.00000000E+00
  0.00000000E+00  0.00000000E+00  0.00000000E+00
  0.00000000E+00  0.00000000E+00  0.00000000E+00
  0.00000000E+00  0.00000000E+00  0.00000000E+00
  0.00000000E+00  0.00000000E+00  0.00000000E+00
  0.00000000E+00  0.00000000E+00  0.00000000E+00
  0.00000000E+00  0.00000000E+00  0.00000000E+00
  0.00000000E+00  0.00000000E+00  0.00000000E+00
  0.00000000E+00  0.00000000E+00  0.00000000E+00
  0.00000000E+00  0.00000000E+00  0.00000000E+00
  0.00000000E+00  0.00000000E+00  0.00000000E+00
  0.00000000E+00  0.00000000E+00  0.00000000E+00
\end{verbatim}
\end{footnotesize}

\subsubsection*{I-43d -- 50~GPa}
\begin{footnotesize}
\begin{verbatim}
Al4C3 I-43d polymorph (findsym-output, f
   1.00000000000000     
     6.0865525100003293    0.0000000000000000    0.0000000000000000
     0.0000000000000000    6.0865525100003293    0.0000000000000000
     0.0000000000000000    0.0000000000000000    6.0865525100003293
   Al   C 
    16    12
Direct
  0.3135103407318738  0.3135103407318738  0.3135103407318738
  0.1864896592681262  0.6864896592681262  0.8135103407318738
  0.6864896592681262  0.8135103407318738  0.1864896592681262
  0.8135103407318738  0.1864896592681262  0.6864896592681262
  0.5635103407318738  0.5635103407318738  0.5635103407318738
  0.9364896592681262  0.4364896592681262  0.0635103407318738
  0.0635103407318738  0.9364896592681262  0.4364896592681262
  0.4364896592681262  0.0635103407318738  0.9364896592681262
  0.8135103407318738  0.8135103407318738  0.8135103407318738
  0.6864896592681262  0.1864896592681262  0.3135103407318738
  0.1864896592681262  0.3135103407318738  0.6864896592681262
  0.3135103407318738  0.6864896592681262  0.1864896592681262
  0.0635103407318738  0.0635103407318738  0.0635103407318738
  0.4364896592681262  0.9364896592681262  0.5635103407318738
  0.5635103407318738  0.4364896592681262  0.9364896592681262
  0.9364896592681262  0.5635103407318738  0.4364896592681262
  0.3750000000000000  0.0000000000000000  0.2500000000000000
  0.1250000000000000  0.0000000000000000  0.7500000000000000
  0.6250000000000000  0.5000000000000000  0.2500000000000000
  0.8750000000000000  0.5000000000000000  0.7500000000000000
  0.2500000000000000  0.3750000000000000  0.0000000000000000
  0.7500000000000000  0.1250000000000000  0.0000000000000000
  0.2500000000000000  0.6250000000000000  0.5000000000000000
  0.7500000000000000  0.8750000000000000  0.5000000000000000
  0.0000000000000000  0.2500000000000000  0.3750000000000000
  0.0000000000000000  0.7500000000000000  0.1250000000000000
  0.5000000000000000  0.2500000000000000  0.6250000000000000
  0.5000000000000000  0.7500000000000000  0.8750000000000000
 
  0.00000000E+00  0.00000000E+00  0.00000000E+00
  0.00000000E+00  0.00000000E+00  0.00000000E+00
  0.00000000E+00  0.00000000E+00  0.00000000E+00
  0.00000000E+00  0.00000000E+00  0.00000000E+00
  0.00000000E+00  0.00000000E+00  0.00000000E+00
  0.00000000E+00  0.00000000E+00  0.00000000E+00
  0.00000000E+00  0.00000000E+00  0.00000000E+00
  0.00000000E+00  0.00000000E+00  0.00000000E+00
  0.00000000E+00  0.00000000E+00  0.00000000E+00
  0.00000000E+00  0.00000000E+00  0.00000000E+00
  0.00000000E+00  0.00000000E+00  0.00000000E+00
  0.00000000E+00  0.00000000E+00  0.00000000E+00
  0.00000000E+00  0.00000000E+00  0.00000000E+00
  0.00000000E+00  0.00000000E+00  0.00000000E+00
  0.00000000E+00  0.00000000E+00  0.00000000E+00
  0.00000000E+00  0.00000000E+00  0.00000000E+00
  0.00000000E+00  0.00000000E+00  0.00000000E+00
  0.00000000E+00  0.00000000E+00  0.00000000E+00
  0.00000000E+00  0.00000000E+00  0.00000000E+00
  0.00000000E+00  0.00000000E+00  0.00000000E+00
  0.00000000E+00  0.00000000E+00  0.00000000E+00
  0.00000000E+00  0.00000000E+00  0.00000000E+00
  0.00000000E+00  0.00000000E+00  0.00000000E+00
  0.00000000E+00  0.00000000E+00  0.00000000E+00
  0.00000000E+00  0.00000000E+00  0.00000000E+00
  0.00000000E+00  0.00000000E+00  0.00000000E+00
  0.00000000E+00  0.00000000E+00  0.00000000E+00
  0.00000000E+00  0.00000000E+00  0.00000000E+00
\end{verbatim}
\end{footnotesize}



\section{État de la campagne}
\label{sec:status}

\begin{center}
\begin{tabular}{p{4.3cm}p{2.4cm}p{4cm}p{3cm}}
\toprule
Polymorphe & Points préparés & Amorçage & État \\
\midrule
Pnma-VII, Pbam, Pnma, R-3m, Pnma-III, anti-CaFe$_2$O$_4$, I-43d
  & 11/11 (0--100~GPa) & WAVECAR PBE, même pression & fichiers prêts, non soumis \\
Fd-3m
  & 11/11 (0--100~GPa) & WAVECAR PBE à $P=0$ seulement, chaîne ascendante ensuite
  & \textbf{soumis} (job SLURM 58461) \\
\bottomrule
\end{tabular}
\end{center}

Les fichiers des sept premiers polymorphes sont générés dans \texttt{hp\_curves\_r2SCAN/}
(\texttt{gen\_hp\_runs\_r2scan.py}) mais pas encore soumis à SLURM. La chaîne Fd-3m a été
soumise séparément (\texttt{gen\_hp\_run\_fd3m\_r2scan.py}), le point $P=0$ étant le seul à
disposer d'un \texttt{WAVECAR} PBE convergé au moment du lancement (le scan PBE de Fd-3m,
lancé en chaîne ascendante unique lui aussi, est toujours en cours -- job SLURM 58460).

\section{Prochaines étapes}

\begin{enumerate}[label=(\roman*)]
  \item Soumettre les sept chaînes r2SCAN restantes (\texttt{hp\_curves\_r2SCAN/submit\_chain.sh}),
  après revue.
  \item Une fois les points r2SCAN convergés, extraire $E$, $V$, $H$
  (\texttt{hp\_curves\_r2SCAN/extract\_hp.py}, même format que le scan PBE) et construire les
  diagrammes $V(P)$, $E(P)$, $H(P)$ r2SCAN, superposés aux courbes PBE du rapport principal.
  \item Comparer les pressions de transition PBE (26 et 52~GPa) aux valeurs r2SCAN
  correspondantes, avec la même analyse de précision ($\delta P_t \approx \sqrt2\,\delta H /
  |\Delta V(P_t)|$) que dans le rapport principal.
  \item Vérifier en particulier si r2SCAN lève ou confirme la quasi-dégénérescence
  Pbam~$\leftrightarrow$~Pnma-III observée en PBE au-delà de 70~GPa.
\end{enumerate}

\appendix
\section{Références du dépôt}

Scripts de génération : \texttt{hp\_curves/gen\_hp\_runs\_r2scan.py} (sept polymorphes convergés
en PBE), \texttt{hp\_curves/gen\_hp\_run\_fd3m\_r2scan.py} (Fd-3m, chaîne ascendante). Données
d'entrée et de sortie dans \texttt{hp\_curves\_r2SCAN/Al4C3\_<polymorphe>/P<pression>/}, script
d'extraction \texttt{hp\_curves\_r2SCAN/extract\_hp.py}. Scan PBE de référence :
\texttt{rapport\_activite\_Al4C3\_HP\_scan.tex} (commit \texttt{effe450}).

\end{document}
